% Configuration du style LaTeX
% Fichier: config/style.tex

% ============================================================================
% STYLES PERSONNALISÉS POUR LE RAPPORT DE SÉCURITÉ
% ============================================================================

% ----------------------------------------------------------------------------
% CONFIGURATION DES POLICES
% ----------------------------------------------------------------------------

% Police principale
\renewcommand{\familydefault}{\sfdefault}

% Tailles de police personnalisées
\newcommand{\titlestyle}{\fontsize{24}{28}\selectfont\bfseries}
\newcommand{\sectionstyle}{\fontsize{16}{20}\selectfont\bfseries}
\newcommand{\subsectionstyle}{\fontsize{14}{18}\selectfont\bfseries}
\newcommand{\subsubsectionstyle}{\fontsize{12}{16}\selectfont\bfseries}
\newcommand{\bodystyle}{\fontsize{11}{13}\selectfont}
\newcommand{\footnotestyle}{\fontsize{9}{11}\selectfont}

% Application des styles de police
\AtBeginDocument{
    \bodystyle
}

% ----------------------------------------------------------------------------
% CONFIGURATION DES COULEURS THÉMATIQUES
% ----------------------------------------------------------------------------

% Palette de couleurs pour la sécurité
\definecolor{primarycolor}{RGB}{0, 123, 255}      % Bleu principal
\definecolor{secondarycolor}{RGB}{108, 117, 125}  % Gris secondaire
\definecolor{successcolor}{RGB}{40, 167, 69}      % Vert succès
\definecolor{dangercolor}{RGB}{220, 53, 69}       % Rouge danger
\definecolor{warningcolor}{RGB}{255, 193, 7}      % Jaune avertissement
\definecolor{infocolor}{RGB}{23, 162, 184}        % Cyan info
\definecolor{lightcolor}{RGB}{248, 249, 250}      % Gris clair
\definecolor{darkcolor}{RGB}{52, 58, 64}          % Gris foncé

% Couleurs pour les niveaux de sévérité
\definecolor{severitycritical}{RGB}{220, 53, 69}   % Rouge critique
\definecolor{severityhigh}{RGB}{253, 126, 20}      % Orange élevé
\definecolor{severitymedium}{RGB}{255, 193, 7}     % Jaune moyen
\definecolor{severitylow}{RGB}{40, 167, 69}        % Vert bas
\definecolor{severityinfo}{RGB}{23, 162, 184}      % Cyan information

% Couleurs pour les catégories OWASP
\definecolor{owaspa01}{RGB}{220, 53, 69}    % A01: Broken Access Control
\definecolor{owaspa02}{RGB}{253, 126, 20}   % A02: Cryptographic Failures
\definecolor{owaspa03}{RGB}{255, 193, 7}    % A03: Injection
\definecolor{owaspa04}{RGB}{40, 167, 69}    % A04: Insecure Design
\definecolor{owaspa05}{RGB}{23, 162, 184}   % A05: Security Misconfiguration
\definecolor{owaspa06}{RGB}{111, 66, 193}   % A06: Vulnerable Components
\definecolor{owaspa07}{RGB}{220, 53, 150}   % A07: Auth Failures
\definecolor{owaspa08}{RGB}{253, 126, 100}  % A08: Integrity Failures
\definecolor{owaspa09}{RGB}{255, 193, 50}   % A09: Logging Failures
\definecolor{owaspa10}{RGB}{40, 167, 120}   % A10: SSRF

% ----------------------------------------------------------------------------
% CONFIGURATION DES MARGES ET ESPACEMENTS
% ----------------------------------------------------------------------------

% Marges personnalisées
\setlength{\parindent}{0pt}
\setlength{\parskip}{6pt}
\setlength{\textheight}{23cm}
\setlength{\textwidth}{16cm}
\setlength{\headheight}{15pt}
\setlength{\headsep}{25pt}
\setlength{\footskip}{30pt}

% Espacement des listes
\setlist[itemize]{topsep=4pt, partopsep=0pt, parsep=2pt, itemsep=2pt, leftmargin=*}
\setlist[enumerate]{topsep=4pt, partopsep=0pt, parsep=2pt, itemsep=2pt, leftmargin=*}

% ----------------------------------------------------------------------------
% STYLES POUR LES ÉLÉMENTS DE SÉCURITÉ
% ----------------------------------------------------------------------------

% Style pour les alertes de sécurité
\newcommand{\securityalert}[2]{%
    \begin{tcolorbox}[
        colback=#1!5!white,
        colframe=#1!75!black,
        title=\textbf{Alerte de Sécurité},
        fonttitle=\bfseries,
        arc=0mm,
        boxrule=0.5pt,
        breakable
    ]
    #2
    \end{tcolorbox}
}

% Style pour les notes importantes
\newcommand{\importantnote}[1]{%
    \begin{tcolorbox}[
        colback=warningcolor!10!white,
        colframe=warningcolor!75!black,
        title=\textbf{Note Importante},
        fonttitle=\bfseries\small,
        before skip=10pt,
        after skip=10pt
    ]
    #1
    \end{tcolorbox}
}

% Style pour les conseils pratiques
\newcommand{\practicetip}[1]{%
    \begin{tcolorbox}[
        colback=successcolor!10!white,
        colframe=successcolor!75!black,
        title=\textbf{Conseil Pratique},
        fonttitle=\bfseries\small,
        before skip=8pt,
        after skip=8pt
    ]
    #1
    \end{tcolorbox}
}

% ----------------------------------------------------------------------------
% STYLES POUR LES VULNÉRABILITÉS
% ----------------------------------------------------------------------------

% Style pour les niveaux de difficulté
\newcommand{\difficultystar}[1]{%
    \textcolor{warningcolor}{\textbf{#1}}
}

% Style pour les catégories OWASP
\newcommand{\owaspcategory}[1]{%
    \ifstrequal{#1}{A01}{\textcolor{owaspa01}{\textbf{#1}}}{%
    \ifstrequal{#1}{A02}{\textcolor{owaspa02}{\textbf{#1}}}{%
    \ifstrequal{#1}{A03}{\textcolor{owaspa03}{\textbf{#1}}}{%
    \ifstrequal{#1}{A04}{\textcolor{owaspa04}{\textbf{#1}}}{%
    \ifstrequal{#1}{A05}{\textcolor{owaspa05}{\textbf{#1}}}{%
    \ifstrequal{#1}{A06}{\textcolor{owaspa06}{\textbf{#1}}}{%
    \ifstrequal{#1}{A07}{\textcolor{owaspa07}{\textbf{#1}}}{%
    \ifstrequal{#1}{A08}{\textcolor{owaspa08}{\textbf{#1}}}{%
    \ifstrequal{#1}{A09}{\textcolor{owaspa09}{\textbf{#1}}}{%
    \ifstrequal{#1}{A10}{\textcolor{owaspa10}{\textbf{#1}}}{%
    \textbf{#1}}}}}}}}}}}
}

% Style pour les niveaux de sévérité
\newcommand{\severitylevel}[1]{%
    \ifstrequal{#1}{Critique}{\textcolor{severitycritical}{\textbf{#1}}}{%
    \ifstrequal{#1}{Élevée}{\textcolor{severityhigh}{\textbf{#1}}}{%
    \ifstrequal{#1}{Moyenne}{\textcolor{severitymedium}{\textbf{#1}}}{%
    \ifstrequal{#1}{Basse}{\textcolor{severitylow}{\textbf{#1}}}{%
    \ifstrequal{#1}{Information}{\textcolor{severityinfo}{\textbf{#1}}}{%
    \textbf{#1}}}}}}
}

% ----------------------------------------------------------------------------
% STYLES POUR LES TABLEAUX
% ----------------------------------------------------------------------------

% Style pour les tableaux de vulnérabilités
\newenvironment{vulnerabilitytable}{%
    \begin{table}[H]
    \centering
    \begin{tabularx}{\textwidth}{>{\raggedright\arraybackslash}X c c c c}
    \toprule
    \textbf{Vulnérabilité} & \textbf{Catégorie} & \textbf{Difficulté} & \textbf{Sévérité} & \textbf{Page} \\
    \midrule
}{%
    \bottomrule
    \end{tabularx}
    \end{table}
}

% Style pour les tableaux de recommandations
\newenvironment{recommendationtable}{%
    \begin{table}[H]
    \centering
    \begin{tabularx}{\textwidth}{>{\raggedright\arraybackslash}X c >{\raggedright\arraybackslash}X}
    \toprule
    \textbf{Recommandation} & \textbf{Priorité} & \textbf{Description} \\
    \midrule
}{%
    \bottomrule
    \end{tabularx}
    \end{table}
}

% ----------------------------------------------------------------------------
% STYLES POUR LES FIGURES
% ----------------------------------------------------------------------------

% Style pour les figures de démonstration
\newcommand{\demoimage}[3]{%
    \begin{figure}[H]
    \centering
    \includegraphics[width=0.8\textwidth]{#1}
    \caption{#2}
    \label{fig:#3}
    \end{figure}
}

% Style pour les figures avec annotations
\newcommand{\annotatedimage}[4]{%
    \begin{figure}[H]
    \centering
    \begin{tikzpicture}
        \node[anchor=south west,inner sep=0] (image) at (0,0) {\includegraphics[width=#1]{#2}};
        \begin{scope}[x={(image.south east)},y={(image.north west)}]
            #4
        \end{scope}
    \end{tikzpicture}
    \caption{#3}
    \end{figure}
}

% ----------------------------------------------------------------------------
% STYLES POUR LES EN-TÊTES ET PIEDS DE PAGE
% ----------------------------------------------------------------------------

% En-tête personnalisé pour les chapitres
\fancypagestyle{chapterstyle}{%
    \fancyhf{}
    \fancyhead[L]{\textcolor{primarycolor}{\nouppercase{\leftmark}}}
    \fancyhead[R]{\textcolor{secondarycolor}{\thepage}}
    \renewcommand{\headrulewidth}{0.4pt}
    \renewcommand{\footrulewidth}{0pt}
}

% Pied de page pour les annexes
\fancypagestyle{appendixstyle}{%
    \fancyhf{}
    \fancyhead[L]{\textcolor{secondarycolor}{Annexe \thechapter}}
    \fancyhead[R]{\textcolor{secondarycolor}{\thepage}}
    \fancyfoot[C]{\textcolor{secondarycolor}{\small Document confidentiel - Distribution limitée}}
    \renewcommand{\headrulewidth}{0.2pt}
    \renewcommand{\footrulewidth}{0.2pt}
}

% ----------------------------------------------------------------------------
% STYLES POUR LES SECTIONS SPÉCIALES
% ----------------------------------------------------------------------------

% Style pour les sections de méthodologie
\newenvironment{methodologysection}{%
    \sectionstyle
    \color{infocolor}
    \begin{tcolorbox}[
        colback=infocolor!5!white,
        colframe=infocolor!75!black,
        before skip=15pt,
        after skip=15pt,
        breakable
    ]
}{%
    \end{tcolorbox}
}

% Style pour les sections de recommandations
\newenvironment{recommendationsection}{%
    \sectionstyle
    \color{successcolor}
    \begin{tcolorbox}[
        colback=successcolor!5!white,
        colframe=successcolor!75!black,
        before skip=15pt,
        after skip=15pt,
        breakable
    ]
}{%
    \end{tcolorbox}
}

% Style pour les sections de vulnérabilités
\newenvironment{vulnerabilitysection}{%
    \sectionstyle
    \color{dangercolor}
    \begin{tcolorbox}[
        colback=dangercolor!5!white,
        colframe=dangercolor!75!black,
        before skip=15pt,
        after skip=15pt,
        breakable
    ]
}{%
    \end{tcolorbox}
}

% ----------------------------------------------------------------------------
% CONFIGURATION DES HYPERRÉFÉRENCES
% ----------------------------------------------------------------------------

% Style pour les liens internes
\newcommand{\internallink}[2]{%
    \hyperref[#1]{\textcolor{primarycolor}{#2}}
}

% Style pour les liens externes
\newcommand{\externallink}[2]{%
    \href{#1}{\textcolor{infocolor}{#2}}
}

% Style pour les références aux figures
\newcommand{\figureref}[1]{%
    \hyperref[fig:#1]{Figure~\ref*{fig:#1}}
}

% Style pour les références aux tableaux
\newcommand{\tableref}[1]{%
    \hyperref[tab:#1]{Tableau~\ref*{tab:#1}}
}

% ----------------------------------------------------------------------------
% STYLES POUR LA PAGE DE TITRE
% ----------------------------------------------------------------------------

% Commande pour le titre principal
\newcommand{\maintitle}[1]{%
    {\titlestyle\color{primarycolor}#1}
}

% Commande pour le sous-titre
\newcommand{\subtitle}[1]{%
    {\sectionstyle\color{secondarycolor}#1}
}

% Commande pour les informations de projet
\newcommand{\projectinfo}[2]{%
    \begin{tabular}{p{0.3\textwidth} p{0.6\textwidth}}
        \textbf{#1:} & #2 \\
    \end{tabular}
}

% ----------------------------------------------------------------------------
% STYLES POUR LES ÉLÉMENTS INTERACTIFS
% ----------------------------------------------------------------------------

% Style pour les boutons (simulés)
\newcommand{\securitybutton}[2]{%
    \tikz[baseline=(char.base)]{
        \node[shape=rectangle, rounded corners=3pt, draw=primarycolor, fill=primarycolor!10, 
              inner sep=4pt, minimum height=20pt] (char) {\textcolor{primarycolor}{\textbf{#1}}};
    }
    \quad #2
}

% Style pour les indicateurs d'état
\newcommand{\statusindicator}[2]{%
    \tikz[baseline=(char.base)]{
        \node[shape=circle, draw=#1, fill=#1!20, inner sep=2pt] (char) {};
    }
    \quad #2
}

% ----------------------------------------------------------------------------
% CONFIGURATION FINALE
% ----------------------------------------------------------------------------

% Application des styles de page
\AtBeginDocument{
    \pagestyle{fancy}
}

% Configuration des sauts de page
\widowpenalty=10000
\clubpenalty=10000
\brokenpenalty=10000

% Fin du fichier style.tex